


% \documentclass[preprint,12pt,authoryear]{elsarticle}
%% Use the option review to obtain double line spacing
%% \documentclass[authoryear,preprint,review,12pt]{elsarticle}

%% Use the options 1p,twocolumn; 3p; 3p,twocolumn; 5p; or 5p,twocolumn
%% Use the options 1p,twocolumn; 3p; 3p,twocolumn; 5p; or 5p,twocolumn
%% for a journal layout:
%% \documentclass[final,1p,times,authoryear]{elsarticle}
% \documentclass[final,1p,times,twocolumn,authoryear]{elsarticle}
% \documentclass[final,3p,times,authoryear]{elsarticle}
%% \documentclass[final,3p,times,twocolumn,authoryear]{elsarticle}
%% \documentclass[final,5p,times,authoryear]{elsarticle}
\documentclass[final,5p,times,twocolumn,authoryear]{elsarticle}

%% For including figures, graphicx.sty has been loaded in
%% elsarticle.cls. If you prefer to use the old commands
%% please give \usepackage{epsfig}

%% The amssymb package provides various useful mathematical symbols
\usepackage{amssymb}
%% The amsmath package provides various useful equation environments.
\usepackage{amsmath}
%% The amsthm package provides extended theorem environments
%% \usepackage{amsthm}

%% The lineno packages adds line numbers. 
\usepackage{lineno}
% Start line numbering with
%% \begin{linenumbers}, end it with \end{linenumbers}. Or switch it on
%% for the whole article with \linenumbers.


\usepackage{kotex}
\usepackage{booktabs}
\usepackage{supertabular}
\usepackage{afterpage}
\usepackage{graphicx}
\graphicspath{{./images/}} 
% hyperlink related 
\usepackage{url}       % URL 줄바꿈을 더 깔끔하게 처리
\usepackage{hyperref}  % 클릭 가능한 링크 생성


% 캡션 정렬.
\usepackage{caption}
\captionsetup[figure]{labelsep=period}

\journal{Applied Energy}

\begin{document}

\begin{frontmatter}




%% Title, authors and addresses

%% use the tnoteref command within \title for footnotes;
%% use the tnotetext command for theassociated footnote;
%% use the fnref command within \author or \affiliation for footnotes;
%% use the fntext command for theassociated footnote;
%% use the corref command within \author for corresponding author footnotes;
%% use the cortext command for theassociated footnote;
%% use the ead command for the email address,
%% and the form \ead[url] for the home page:

%% \title{Title\tnoteref{label1}}
%% \tnotetext[label1]{}
%% \author{Name\corref{cor1}\fnref{label2}}
%% \ead{email address}
%% \ead[url]{home page}
%% \fntext[label2]{}
%% \cortext[cor1]{}
%% \affiliation{organization={},
%%            addressline={}, 
%%            city={},
%%            postcode={}, 
%%            state={},
%%            country={}}
%% \fntext[label3]{}


\title{Contextually Extended iTransformer (CEiT) for Query-Based Virtual Sensing in Dynamic Spatial Sensor Networks} %% Article title

%% 추가 후보 
% Contextually Extended iTransformer for Robust Virtual Sensing with Missing and Variable Sensors

% Contextually Extended iTransformer for Virtual Sensing of Indoor Air Quality in Livestock Buildings



%% 저자 및 소속 설정
\author[a]{Wonjun Choi\fnref{equal}}

\author[b]{Sangwon Lee\fnref{equal}}

\author[c]{Hakjong Shin}

\author[d]{Younghoon Kwak\corref{cor1}}
\ead{ikyh2@uos.ac.kr} %% 마지막 저자의 이메일

%% 하단 주석 설정
\fntext[equal]{These authors contributed equally to this work.}
\cortext[cor1]{Corresponding author}

%% 소속 정보 (이전과 동일)
\affiliation[a]{
            organization={School of Architecture, Chonnam National University}, 
            addressline={77 Yongbong-ro, Buk-gu}, 
            city={Gwangju}, 
            postcode={61186},
            country={South Korea}
            }

\affiliation[b]{organization={Org B},
            addressline={Address B},
            city={City B},
            postcode={67890},
            country={Country B}}

\affiliation[c]{organization={Org C},
            addressline={Address C},
            city={City C},
            postcode={67890},
            country={Country C}}

\affiliation[d]{organization={Org D},
            addressline={Address D},
            city={City D},
            postcode={67890},
            country={Country D}}
\begin{abstract}
%% Text of abstract
abstract 
for real 



aaa 한글가능?





\end{abstract}

%%Graphical abstract
% \begin{graphicalabstract}
%\includegraphics{grabs}
% \end{graphicalabstract}

%%Research highlights
% \begin{highlights}
% \item Research highlight 1
% \item Research highlight 2
% \end{highlights}

%% Keywords
\begin{keyword}

Virtual sensor \sep Deep learning \sep Transformer \sep Smart farm \sep  Spatio-temporal interpolation \sep Interpretable AI 

%% keywords here, in the form: keyword \sep keyword

%% PACS codes here, in the form: \PACS code \sep code

%% MSC codes here, in the form: \MSC code \sep code
%% or \MSC[2008] code \sep code (2000 is the default)

\end{keyword}

\end{frontmatter}



%% Add \usepackage{lineno} before \begin{document} and uncomment 
%% following line to enable line numbers
%% \linenumbers




\renewcommand{\figurename}{Fig.} % 메인 텍스트용 figure 표시 설정

%%%%%%%%%%%%%%%%%%%%% MAIN TEXT 
%% Use \section commands to start a section
% \begin{linenumbers}
% \modulolinenumbers[2] % 5, 10, 15... 행에만 번호 표시

\section{Introduction}
\label{sec:intro}

\afterpage{\afterpage{
\begin{table}[t!]
\section*{Nomenclature}
\begin{tabular}{@{} p{0.15\columnwidth} p{0.8\columnwidth}}
\multicolumn{2}{@{}l}{\textit{Latin Symbols}} \\
$B_{\text{train}}$ & Training batch size \\
$d_{\text{model}}$ & Hidden size of model \\
lr & Adam learning rate \\
$N_{\text{head}}$ & Number of MHA heads \\
$N_{\text{layer}}$ & Number of head layers \\
$N_{OL}$ & Number of layers in output head \\
$N_q$ & Number of query tokens \\
$N_s$ & Number of input sensors \\
$N_{VP}$ & Number of layers in value projector \\
$p$ & 3D spatial coordinate vector $(x, y, z)$ \\
$p_{\text{dropout}}$ & Dropout rate \\
$p_{\text{self}}$ & Self-Reference Training probability \\
$W$ & Random matrix for Gaussian Fourier Features \\
\vspace{0.1cm} \\
\multicolumn{2}{@{}l}{\textit{Greek Symbols}} \\
$\gamma(p)$ & Position embedding function \\
$\sigma$ & Standard deviation for Gaussian Fourier Features \\
\end{tabular}
\end{table}
}}



%% Labels are used to cross-reference an item using \ref command.

\subsection{Background: 가상 센싱의 필요성과 기존 연구}
\label{subsec1}

현대 산업 시스템은 물리 세계의 상태를 디지털 공간에서 실시간으로 관측·해석·제어하는 Cyber-Physical Systems (CPS) \cite{lee2008isorc_cps, rajkumar_cyber-physical_2010}로 진화하고 있으며, 센서 기반 시계열 데이터는 상태 추정과 제어·의사결정의 기반이 된다. IoT 패러다임의 확산으로 센서·액추에이터 네트워크가 대규모로 구축되고, 데이터 기반 운영 최적화가 생산성과 경제성에 직접 연결되는 구조가 일반화되고 있다 \cite{gubbi_internet_2013}. 그러나 데이터 기반 운영의 전제인 데이터 무결성(data integrity)은 현장에서 쉽게 훼손된다. 현실에서 실제 센서는 고장, 드리프트, 통신 장애, 전원 문제에 취약하고, 고온·고압·부식성 가스 등 극한 조건에서는 설치 자체가 제한되거나 장비 수명이 급격히 단축된다 \cite{isermann_model-based_2005}. 이는 결과적으로 결측과 이상치를 빈번히 발생시키며, 이는 진단·제어의 안정성을 훼손한다 \cite{agbo_missing_2022}. 
% 결과적으로, 필요한 수준의 공간 커버리지와 시간 연속성을 물리 센서만으로 달성하기 어렵다.


특히 넓고 복잡한 실내 공간의 환경을 정밀하게 모니터링하기 위해서는 고밀도 센서 배치가 요구된다. 그러나 이러한 요구를 충족하기 위해 물리적 센서를 고밀도로 설치하는 기존의 접근 방식은 심각한 현실적 장벽에 부딪힌다. 첫째, 비용의 비효율성이다. 넓은 공간을 사각지대 없이 커버하기 위해 수십, 수백 개의 고정밀 센서를 설치하고 네트워크를 구축하는 비용은 막대하다. 둘째, 유지보수의 난해함이다. 고온 다습하거나 부식성 가스가 상시 존재하는 가혹한 조건(harsh environment)에서 전자 장비인 센서는 빈번하게 오작동하거나 수명이 단축되며, 잦은 교체와 캘리브레이션 작업은 운영자에게 큰 부담을 지운다. 셋째, 데이터의 불연속성이다. 센서의 고장이나 통신 장애로 인한 데이터 결측(missing data)은 자동화된 제어 시스템의 오작동을 유발하여 시스템 전체의 안정성을 위협할 수 있다.

현대의 스마트 농업과 산업 환경에서는 온도, 습도, 가스 농도와 같은 핵심 환경 변수를 정밀하게 관리하기 위해 이기종 센서 네트워크에 대한 의존도가 점점 높아지고 있다 \cite{wolfert_big_2017}. 특히 양돈 축사와 같은 대규모 실내 시설에서는 적정 온도와 습도, 이산화탄소와 암모니아 농도 등 공기질을 적정 수준으로 유지하는 것이 가축의 건강과 생산성에 직접적인 영향을 미친다 \cite{costa_ammonia_2017, drummond_effects_1980}. 이러한 맥락에서, 제한된 수의 센서 측정값만으로 공간 전역 혹은 원하는 공간좌표의 환경 상태를 추정하는 공간 보간(spatial interpolation) 또는 가상 센싱(virtual sensing) 접근 방식이 효율적이고 현실적인 대안으로 부상하고 있다 \cite{ploennigs_virtual_2011}. 해당 기술의 본질적인 난제는 단순히 소수의 센서로 공간을 보간하는 것을 넘어 서로 다른 유형, 위치, 신뢰성을 가진 불규칙 센서 네트워크에서 생성되는 시공간 데이터 스트림의 특성을 단일 모델이 얼마나 잘 융합하고 예측하느냐에 있다. 이는 공간적으로 인접한 센서 간의 상관성뿐만 아니라 시간 지연을 동반하는 전파·환기 효과, 센서 타입(온도·습도·가스) 간 비대칭적 상호작용, 센서 고장·누락으로 인한 입력 구성 변화가 동시에 발생하고 한 위치의 측정값이 다른 위치의 값을 예측하거나 영향을 미친다. 따라서 기존의 단순한 공간 상관 가정이나 고정된 입력 방식에서 탈피하여, 공간적, 시간적, 그리고 교차 변수 간의 유기적인 종속성을 통합적으로 포착할 수 있는 고도화된 모델링 기법의 개발이 필수적이다.

% 전통적으로 지구통계학 기반 방법(variogram–kriging)은 공간 상관을 이용해 미계측 지점을 추정하는 표준 접근으로 자리 잡아 왔다 [13]. 그러나 실제 환경에서는 비선형·비정상성(국소적 환기, 장애물, 열원 분포 등)이 강해 단일한 상관 구조로 설명이 어려운 경우가 많고, 센서 수가 증가하면 계산 비용이 급격히 커져 대규모 적용이 부담이 된다 [13,14]. 이런 이유로 Kriging 자체의 근사/분할 기반 확장이 제안되어 왔지만, 다변량·이기종 센서 스트림을 유연하게 통합하는 데는 여전히 제약이 남는다 [14].


이러한 필요성에 대응해 다양한 보간 기법이 제안되어 왔다. 전통적으로 Kriging과 같은 통계 기반 기법은 공간 상관(variogram)을 가정하고 미계측 지점을 최적 추정하는 표준 접근으로 사용되어 왔다 \cite{oliver_tutorial_2014}. Kriging은 센서 관측값들 사이의 공간적 상관성을 가정하여 미계측 지점의 값을 최적 추정하지만, 실제 환경에서 센서 간 관계가 복잡하거나 비선형적일 경우 이러한 가정이 성립하지 않아 정확도가 떨어질 수 있다. 또한 고차원 데이터나 다변량 공간 필드(예: 온도와 가스 농도의 동시 추정)로 확장하면 모델링과 파라미터 추정이 급격히 복잡해진다 [CITE]. 또한 관측점 수가 증가할수록 계산 복잡도가 O(n³)로 증가해 대규모·실시간 적용에 제약이 된다 \cite{van_stein_cluster-based_2020}.


한편, 최근에는 CNN과 RNN 계열의 결합, 그래프 신경망(GNN), 트랜스포머(Transformer) 등의 딥러닝 기반 보간·예측 모델이 제안되며 센서 데이터의 공간-시간 패턴 학습에 활용되고 있다 \cite{shi2015convlstm,yu2017stgcn,vaswani2017transformer,zhou2020informer}. 일부 벤치마크에서는 전통적 Kriging 대비 성능 향상이 보고되었지만 \cite{yu2017stgcn,zhou2020informer}, 많은 접근이 ‘현장 조건에서의 입력 유연성’보다는 ‘특정 입력 정의를 전제로 한 아키텍처’에 최적화되어 있어 실제 적용에서 구조적 취약점이 드러난다 \cite{yu2017stgcn,kipf2016gcn}. 예를 들면 센서 간 연결정보를 사전에 알아야 하거나, 입력 형태와 모델 구조가 고정되어 있어 센서 구성 변화나 부분적 센서 데이터 손실에 대응할 수 없는 것이다. 

템포럴 합성곱 신경망(Temporal CNN)혹은 CNN과 RNN(CNN-LSTM) 계열을 결합한 하이브리드 아키텍처는 센서 값을 격자(grid)나 이미지 형태로 재구성해 공간 특징을 추출하고, RNN이나 temporal CNN으로 시간 패턴을 학습한다 \cite{shi2015convlstm}[CITE]. 그러나 이 방식은 불규칙한 센서 네트워크를 다루기 위해 배치를 격자로 강제하는 전처리(리샘플링/패딩)를 요구하고, 센서 추가·제거 및 결측 발생 시 입력 형상이 변해 모델 재구성이나 재학습으로 이어지기 쉽다. 또한 공간 모듈(CNN)과 시간 모듈(RNN)이 분리된 모듈식 설계는 공간·시간 상호작용을 상위 계층에서 자연스럽게 공동 모델링하는 데 구조적 한계를 가질 수 있다 [CITE]. 이를 보완하기 위한 복잡한 게이팅 구조가 제안되었지만, 이는 뉴럴 네트워크 복잡도를 불필요하게 증가시키며, 새로운 도메인이나 센서 구성 변화에 대한 막대한 재설계 비용이 실용성을 제한한다 \cite{jing2019mlclstm}.

합성곱 신경망(CNN)을 사용하는 하이브리드 접근 역시 유사한 제약을 갖는다. CNN은 센서 그리드를 이미지와 유사한 형태로 간주하여 공간적 특징을 추출하고, RNN이나 temporal CNN이 시간 축을 처리하는 방식이 대표적이다 \cite{shi2015convlstm}. 이러한 모듈식 구조는 특정 상황에서는 효과적일 수 있으나, 공간 모듈(CNN/GNN)과 시간 모듈(RNN/CNN)이 분리되어 설계되므로, 상위 계층에서 공간·시간 특성을 공동으로 모델링하고 자연스럽게 융합하는 데 구조적 한계를 가진다 \cite{shi2015convlstm}. 이를 보완하기 위해 다중 게이팅 메커니즘을 갖춘 상관 LSTM 등 보다 정교한 아키텍처가 제안되었지만 \cite{jing2019mlclstm}, 네트워크 구조가 지나치게 복잡해지고 새로운 센서 구성이나 도메인으로 확장할 때 재설계 비용이 커지는 문제가 있다 \cite{jing2019mlclstm}. 요약하면, 많은 기존 시공간 모델은 센서 연결성, 좌표계, 격자 구조와 같은 수동 구성 요소나 특정 시나리오를 전제로 한 특수 설계에 의존하며, 이기종 센서 유형과 동적으로 변화하는 센서 네트워크에 직면했을 때 번거로운 재설정을 요구한다.


이에 대한 대안으로 등장한 그래프 신경망(GNN)\cite{scarselli2009tcnn_gnn} 기반 접근법은 센서를 노드로, 센서 간 영향을 엣지로 표현해 불규칙 배치를 직접 다룰 수 있다는 장점이 있다. GNN과 순환/어텐션 모듈을 결합한 하이브리드 모델은 다양한 시공간 예측에서 우수한 성능을 보고했지만 [CITE], 다수 방법이 인접행렬(연결성)을 사전에 정의하거나 별도로 학습해야 한다는 전제를 가진다. 실제 현장에서는 센서 토폴로지가 변동하고 연결성이 암묵적이거나 불완전한 경우가 많아, 그래프 구성·갱신 자체가 병목이 되며 모델 재학습 또는 구조 변경이 요구될 수 있다.

한편, 순차 데이터 모델링에서 탁월한 성과를 보인 Transformer 계열 모델은 장거리 종속성을 포착하는 능력으로 주목받고 있다\cite{vaswani2017transformer,zhou2020informer}. Transformer는 자기어텐션(self-attention) 메커니즘을 통해 서로 다른 센서 종류(피처)와 시간적 종속성 관계를 명시적으로 학습할 수 있으며, 순차 데이터 내 장거리 종속성과 복잡한 관계를 자동으로 포착할 수 있다 \cite{vaswani2017transformer}. 다양한 시계열 예측 문제에서 우수한 성능을 보여 왔지만, 앞서 설명한 다른 딥러닝 아키텍처(GNN, CNN-LSTM)와 마찬가지로 입력 구조의 유연성 문제에서 자유롭지 못하다. 전통적인 Transformer는 각 시점(time step)을 토큰으로 처리하는 방식으로, 입력 시퀀스 길이가 고정되어야 하거나 토큰 순서에 의존적이며, 센서 개수가 변경되거나 입력 구성이 달라질 경우 모델 재구성이 필요하다. 또한 그래프 인접 행렬이나 격자 구조와 같은 명시적인 공간 구조 정의는 필요하지 않지만, 센서 연결성이나 좌표계 정보를 활용하기 어려워 공간적 관계를 효과적으로 모델링하는 데 한계가 있다 \cite{zhou2020informer,liu2023itransformer}.

% 한편, 순차 데이터 모델링에서 탁월한 성과를 보인 Transformer 계열 모델은 장거리 종속성을 포착하는 능력으로 주목받고 있다 \cite{vaswani2017transformer,zhou2020informer}. 특히 최신 모델인 iTransformer \cite{liu2023itransformer, liu2024iclr_itransformer}는 시점(time step)이 아닌 변수(variable) 자체를 토큰화하는 변수 중심 접근을 통해, 명시적인 그래프 구조 없이도 다변량 종속성을 효과적으로 학습하며 기존 Transformer 계열을 능가하는 성능을 보고했다. 하지만 이러한 최신 아키텍처조차 공간 센서 네트워크 환경에서는 치명적인 한계를 드러낸다. 첫째, iTransformer는 각 센서의 시계열 패턴에 집중할 뿐, 해당 센서의 공간적 위치나 물리적 타입(온도, 습도, 가스 등)과 같은 핵심 메타정보를 명시적으로 인코딩하지 않는다. 물리적으로 인접하거나 상호 비대칭적인 영향을 주고받는 이기종 센서 간의 관계를 파악하기 위해서는 이러한 메타정보가 필수적이다. 둘째, 센서의 고장, 추가, 제거 등으로 입력 집합이 변동되는 상황에서, 고정된 변수 집합을 가정하는 기존 모델들은 구조적 대응이 불가능하거나 성능이 급격히 저하된다. 결과적으로, 센서의 위치 및 타입 정보를 적극적으로 활용하면서도, 동적으로 변화하는 센서 구성에 견고하게 적응할 수 있는 가상 센싱 아키텍처는 여전히 중요한 연구 공백(Research Gap)으로 남아 있다.



최근 제안된 iTransformer는 Transformer의 입력 구조 유연성 문제를 해결하기 위한 새로운 아키텍처로 주목받고 있다 \cite{liu2023itransformer}. 기존의 Transformer가 각 타임스텝의 모든 변수를 토큰으로 처리하는 방식과 달리, iTransformer는 각 변수(센서)를 토큰화하는 변수 중심 토큰화를 채택한다. 즉 각 센서의 전체 과거 시계열을 하나의 벡터로 인코딩하여 토큰을 구성하고, 자기어텐션을 이 변수 토큰들 간에 적용함으로써 센서 간(교차 변수) 상관관계를 학습하고, 동시에 각 토큰 내부의 피드포워드 네트워크가 해당 센서의 시간적 역학을 모델링한다. 이러한 아키텍처는 그래프 인접 행렬이나 격자 구조와 같은 명시적인 공간 구조 정의 없이도 어텐션을 통해 센서 간 관계를 암묵적으로 학습할 수 있어, 입력 구조의 유연성을 크게 향상시켰다. 그 결과, 여러 시계열 포어캐스팅 벤치마크에서 기존 Transformer 및 특화 모델 대비 우수한 성능이 보고되었다 \cite{liu2023itransformer}. 그러나 센서 네트워크를 고려한 공간 보간의 예측 어플리케이션에서는 몇 가지 한계가 있다. 첫째, 변수 토큰화는 각 센서의 시계열을 압축할 뿐 센서의 공간 위치나 센서 타입(예: 온도, 습도, 가스 센서)을 명시적으로 인코딩하지 않아, 물리적 인접성이나 타입 간 비대칭 상호작용 같은 핵심 단서를 모델이 안정적으로 활용하기 어렵다. 예를 들어, 물리적으로 인접한 두 온도 센서는 멀리 떨어진 센서들보다 더 높은 상관관계를 보이는 경향이 있고, 온도·습도·가스 농도와 같은 서로 다른 센서 유형은 유사한 패턴을 공유하기도 하지만 상호 영향은 비대칭적일 수 있다. 둘째, 센서 추가·제거·고장으로 입력 집합이 변동할 때 학습 시점과 다른 변수 집합에 대해 동작과 성능을 일관되게 보장하기 어렵다 \cite{liu2024iclr_itransformer}. 기존 딥러닝 기반 접근법들과 마찬가지로 특정 그래프 또는 그리드 형태의 입력 정의에 의존하거나 정해진 센서 배열에서만 동작하므로, 실제 환경에서 센서 추가·제거 또는 고장으로 인한 데이터 결손 상황에 적용하기 어렵다. 결국 기존 딥러닝 기반 접근은 격자/그래프/고정 변수 집합 등 ‘입력 정의의 고정성’에 기대는 경향이 강하며, 위치·타입 메타정보를 활용하면서도 입력 센서 구성이 동적으로 변해도 견고하게 작동하는 가상 센싱을 위한 아키텍처는 여전히 중요한 연구 공백(research gap)으로 남아 있다.



\subsection{Proposed Architecture: Contextually Extended iTransformer}
\label{subsec2}

본 연구는 위에서 지적한 한계점들—물리 센서망의 비용·운영 제약, 통계 모델의 단순 가정, 기존 딥러닝 모델의 입력 경직성—을 동시에 극복하고, 유연한 가상 센싱 기술을 실현하기 위해 Contextually Extended iTransformer(CEiT) 모델을 제안한다. CEiT는 iTransformer 백본을 기반으로 하면서, 공간 센서 네트워크 환경에서의 한계를 해결하기 위해 두 가지 핵심 기술을 도입한 구조적 확장이다.

첫째, 컨텍스트 임베딩(Contextual Embedding)이다. CEiT는 각 센서 토큰에 대해 센서의 위치 좌표(공간 임베딩)와 센서 유형 식별자(타입 임베딩)를 결합한 학습 가능한 컨텍스트 임베딩을 정의한다. 이를 통해 모델은 물리적으로 가까운 센서일수록 서로 강하게 주의를 기울여야 한다는 사실이나, 공간적으로 일정한 기울기나 전파 지연과 같은 패턴을 데이터로부터 학습할 수 있다. 동시에 센서 타입 정보를 활용함으로써 유사한 변수(예: 온도 계열 센서)를 그룹화하면서도 서로 다른 모달리티를 구분할 수 있다. 즉, 각 변수 토큰이 "어디에서 무엇을 측정하는지"에 대한 정보를 내포하게 함으로써, Transformer의 어텐션 메커니즘이 수동으로 정의된 그래프 없이도 시공간 및 유형 간 종속성을 암묵적으로 발견할 수 있도록 한다.

% REVIEW: 컨텍스트 임베딩의 구성 요소(공간 임베딩 + 타입 임베딩)와 그 효과를 명확히 설명하고 있다. "물리적으로 가까운 센서일수록 서로 강하게 주의를 기울여야 한다"는 점이 공간 임베딩의 핵심 가치를 잘 드러내고 있으며, "공간적으로 일정한 기울기나 전파 지연"과 같은 구체적 패턴 예시가 좋다. 마지막 문장의 "암묵적으로 발견"이라는 표현이 기존 GNN의 명시적 그래프 정의와의 차별점을 잘 강조하고 있다.

둘째, 쿼리 토큰(Query Token)의 도입이다. Vision Transformer (ViT)\cite{dosovitskiy2021iclr_vit}의 설계에서 영감을 받아, CEiT는 Transformer 인코더에 특정 센서 입력에 직접 대응하지 않는 학습 가능한 쿼리 토큰을 도입한다.  ViT에서 CLS(Classification) 토큰이 패치 토큰들로부터 분류에 필요한 정보를 집약하듯, CEiT의 쿼리 토큰은 모든 센서 토큰으로부터 예측에 필요한 정보를 능동적으로 수집한다. 이를 통해 사용자는 특정 공간 좌표(x, y, z) 에서의 온도 값과 같이 임의의 지점에 대한 환경 상태를 질의 형태로 지정할 수 있으며, 쿼리 토큰은 어텐션을 통해 필요한 정보를 집계한다. 중요한 점은, 새로운 예측 대상을 정의할 때마다 모델을 다시 학습할 필요가 없으며, 부분적인 센서 하위 집합만으로도 관심 대상에 집중된 예측을 수행할 수 있어, 센서 고장이나 누락 상황에서도 견고하게 작동한다는 것이다.

이와 같이 CEiT는 iTransformer가 가진 장점(변수 간 의존성 포착, 효율적인 인코더 구조)을 유지하면서, 공간 센서 네트워크 환경에서 치명적이었던 위치·타입 정보 부재 문제를 해결하고, 명시적 그래프 정의 없이도 시공간 상관관계를 학습하는 유연하고 구조적으로 간결한 시공간 예측 프레임워크를 제공한다. 자기어텐션(self-attention) 메커니즘을 통해 센서들 간의 공간 상관관계를 인접행렬 없이도 학습하며, 사전 정의된 그래프 구조에 의존하지 않기 때문에 입력 센서의 일부 결손이나 구성 변화에도 유연하게 대응할 수 있다.


본 연구의 주요 기여는 다음과 같다:

\begin{itemize}
    \item \textbf{입력에 유연하고 강건한 공간 보간 아키텍처 제안:} 그래프 구조 정보 없이도 센서 간 복잡한 관계를 학습하는 CEiT 모델을 개발함으로써, 센서 고장이나 추가/제거 상황에서도 안정적인 예측이 가능하다.
    
    \item \textbf{컨텍스트 임베딩 및 쿼리 토큰을 활용한 가상 센싱:} 센서의 위치좌표 및 종류 메타정보를 임베딩하여 모델 입력에 활용하고, 쿼리 토큰을 통해 임의 위치의 환경 값을 직접 질의하는 새로운 방법을 제시하였다. 
    
    \item \textbf{실측 데이터 기반 성능 검증:} 실제 축사 센서 데이터(온도, 습도, CO₂, NH₃)를 활용해 공간 보간과 센서 결측/고장 복원 시나리오에서 성능과 견고성을 평가했다.
\end{itemize}

이상의 결과로부터, 제안하는 CEiT 모델은 넓은 공간의 환경 모니터링을 위한 가상 센서 솔루션으로서 높은 실용성과 성능을 나타내며, 복잡한 실세계 환경에서의 공간 보간 문제를 효과적으로 해결할 수 있음을 확인하였다.

% REVIEW: Contributions 섹션은 세 가지 주요 기여를 명확하게 나열하고 있으며, 각 기여가 Background에서 제시한 연구 공백을 어떻게 해결하는지 연결되어 있다. 다만, 각 기여의 구체적인 수치적 성과나 정량적 비교 결과는 Results 섹션에서 다루어야 하므로, Introduction에서는 이 정도 수준의 언급이 적절하다.


%% Use \subsubsection, \paragraph, \subparagraph commands to 
%% start 3rd, 4th and 5th level sections.
%% Refer following link for more details.
%% https://en.wikibooks.org/wiki/LaTeX/Document_Structure#Sectioning_commands




%% Use \section commands to start a section
\section{Data}

%% Use \subsection commands to start a subsection.
\subsection{Measurement}
measurement 


%% Use \subsection commands to start a subsection.
\subsection{Preprocessing and feature engineering}

ahahahahahahhaha




%% CEiT Architecture figure
\begin{figure*}[htbp] % 위치 우선순위: 현재 위치(h) -> 위(t) -> 아래(b) -> 별도 페이지(p)
    \centering
    \includegraphics[width=0.95\linewidth]{ceit_architecture.png} % 0.95: 전체너비의 95%
    \caption{CEiT architecture and components detail. (a) Architecture overview, (b) Tokenization from multi modal inputs, (c) Tokenization of queries, and (d) Hierarchical attention masking }
    \label{fig:ceit_architecture}
\end{figure*}

%% Use \section commands to start a section
\section{Contextually Extended iTransformer}

\subsection{모델 아키텍처 개요}


본 연구에서는 축사 내 환경 모니터링을 위한 가상 센서 구현을 목표로, 공간적 보간에 특화된 모델을 제안한다. 제안된 모델은 iTransformer (Inverted Transformer) \cite{liu2024iclr_itransformer} 아키텍처를 기반으로 하되, 시계열 예측이 아닌 동일 시간대의 공간적 보간 작업에 최적화되도록 설계되었다. 핵심 개선점은 센서의 타입 정보와 3차원 공간 좌표를 효과적으로 통합하여, 제한된 수의 물리 센서로부터 임의 위치의 환경 값을 정확하게 추정하는 것이다.
전통적인 Transformer 모델이 서로 다른 타임스텝 사이의 관계를 학습하는 것과 달리, 본 모델은 센서를 독립적인 토큰으로 취급하여 센서 간 공간적 관계를 학습한다. 각 센서는 24시간(48개 타임스텝)의 시계열 데이터를 포함하고 있으며, 이는 일중 변화 패턴과 시간적 상관관계를 포착하는 컨텍스트로 활용된다. 이러한 접근법을 통해 모델은 시공간적 패턴을 동시에 학습하여 높은 정확도의 공간 보간을 달성한다.

% figure 대신 figure* 써야 한다고 함.
\begin{figure*}[tp]%% placement specifier
%% Use \includegraphics command to insert graphic files. Place graphics files in 
%% working directory.
  \centering
  \fbox{\parbox{0.9\textwidth}{\centering \vspace{2cm} FIGURE PLACEHOLDER \vspace{2cm}}}
%% Use \caption command for figure caption and label.
\captionsetup{width=\textwidth}
\caption{모델 설명하는 그림, 각 토큰이 압축된 시계열 + 타입, 위치 임베딩임 표시, 구체적인 트랜스포머 구조는 선택사항}
%% https://en.wikibooks.org/wiki/LaTeX/Importing_Graphics#Importing_external_graphics
\end{figure*}

% REVIEW: The overview effectively contrasts the proposed model with traditional Transformers. However, it is crucial to explicitly distinguish this from the original iTransformer, which also treats variates as tokens. Clarify if the "spatial relationship" learning is an emergent property of the attention mechanism or explicitly enforced. Also, the reliance on a fixed 24-hour (48-step) window should be justified—is this hardcoded, and how does it affect real-time applicability?


\subsection{멀티모달 임베딩 통합}
제안된 모델의 첫 번째 핵심 구성요소는 세 가지 서로 다른 정보원을 통합하는 멀티모달 임베딩 시스템이다. 각 센서 토큰은 측정값, 공간 위치, 그리고 센서 타입이라는 세 가지 독립적인 정보를 포함하고 있으며, 이들은 각각 전용 임베딩 모듈을 통해 처리된 후 합산되어 최종 토큰 표현을 형성한다.

% 모델 아키텍처 그림 필요할듯? 


\subsubsection{값 임베딩 (Value Embedding)}
각 센서의 48개 타임스텝 측정값은 먼저 다층 퍼셉트론(MLP)을 통해 모델의 히든 차원($d_{model}=768$)으로 변환된다. 이 과정에서 원시 시계열 데이터의 시간적 패턴이 압축되어 표현되며, 학습 과정에서 각 센서 타입별 특성에 맞는 표현이 자동으로 학습된다. MLP의 깊이는 하이퍼파라미터로 조정 가능하며, 실험을 통해 3층 구조가 최적임을 확인하였다.
% REVIEW: Specifying "48 time steps" implies a dependency on a fixed input window size. Consider explaining how the model handles missing steps or variable-length inputs if applicable. If strictly 48, state this limitation clearly.


\subsubsection{위치 임베딩 (Position Embedding)}
3차원 공간 정보의 효과적인 인코딩을 위해 Gaussian Fourier Features 기법을 채택하였다. 이 방법은 저차원 공간 좌표를 고차원 특징 공간으로 매핑하여, 모델이 공간적 근접성과 거리 관계를 효과적으로 학습할 수 있도록 한다. 
구체적으로, 3차원 좌표 벡터 $p = (x, y, z)$ 는 먼저 평균 0, 표준편차 $\sigma$ 의 가우시안 분포에서 샘플링된 랜덤 행렬 $W  \in R^{3 \times (d/2)}$ 와 곱해진다. 이후 사인과 코사인 함수를 적용하여 최종 $d$ 차원 위치 임베딩을 생성한다. 수식으로 표현하면 다음과 같다:

\begin{equation}
    \gamma(p) = \left[\sin(2\pi W p), \cos(2\pi W p)\right]
\end{equation}
이 인코딩 방식의 주요 장점은 학습 파라미터 없이도 공간적 유사성을 보존할 수 있다는 점이다. 또한, 외부 기상 센서와 같이 특정 위치가 정의되지 않은 센서들을 위해 학습 가능한 fallback 파라미터를 제공하여, 모든 센서 타입을 통일된 프레임워크 내에서 처리할 수 있도록 하였다.  $\sigma$ 값은 하이퍼파라미터 선택 과정에서 최적화 되어 선택된다.
% REVIEW: Clarify if the random matrix $W$ is fixed (frozen) after initialization or learned. Standard GFF uses fixed random weights, which preserves the distance metric properties. If learned, it becomes a standard position embedding initialized with Fourier features.


\subsubsection{타입 임베딩 (Type Embedding)}
센서 타입 정보는 학습 가능한 임베딩 테이블을 통해 인코딩된다. 온도, 습도, 이산화탄소, 암모니아 등 각 센서 타입은 고유한 $d$차원 벡터로 표현되며, 이는 학습 과정에서 각 타입의 특성과 다른 타입과의 관계를 반영하도록 최적화된다. 이를 통해 모델은 동일한 위치에서도 센서 타입에 따라 다른 예측을 수행할 수 있다.
세 가지 임베딩은 요소별 합산을 통해 통합되어 최종 센서 토큰 표현을 형성한다. 이러한 가법적 통합 방식은 각 정보원의 독립성을 유지하면서도 상호작용을 가능하게 하여, 모델이 복잡한 시공간적 패턴을 학습할 수 있도록 한다.

\subsection{쿼리 기반 가상 센서 예측 메커니즘}

\subsubsection{쿼리 토큰의 개념과 역할}
본 모델의 주요 차별점은 예측 대상인 가상 센서를 "쿼리 토큰(Query Token)"으로 모델링하는 접근법이다. 전통적인 회귀 모델이 고정된 출력 차원을 가지는 것과 달리, 제안된 모델은 임의의 수와 위치의 가상 센서를 동적으로 정의하고 예측할 수 있다. 이는 Transformer의 sequence-to-sequence 특성을 활용하여, 예측하고자 하는 각 위치를 독립적인 쿼리 토큰으로 표현함으로써 달성된다.
쿼리 토큰은 두 가지 정보를 기반으로 구성된다. 첫째, 예측하고자 하는 3차원 공간 좌표 $p=(x, y, z)$로, 이는 가상 센서가 설치될 위치를 나타낸다. 둘째, 센서 타입 정보로, 해당 위치에서 예측하고자 하는 환경 변수(온도, 습도, CO2, 암모니아)를 지정한다. 입력 토큰과 달리 쿼리 토큰은 값 임베딩을 포함하지 않으며, 위치 임베딩과 타입 임베딩의 합으로만 구성된다. 이러한 구조적 차이는 쿼리 토큰이 "질문"의 역할을 수행함을 명확히 한다.
예를 들어, 축사의 중앙 지점 (3500, 1500, 1800)에서 온도를 예측하고자 한다면, 해당 위치의 Gaussian Fourier Features 인코딩과 온도 타입 임베딩을 합산하여 쿼리 토큰을 생성한다. 이 토큰은 Transformer 처리 과정에서 입력 센서들의 정보를 종합하여, 최종적으로 Output MLP를 통해 해당 위치의 온도값 시계열을 출력한다.

% REVIEW: The "Query Token" concept is a strong novelty here, analogous to the [CLS] token or DETR's object queries. It effectively turns the interpolation problem into a query-response task. This paragraph is clear and logical.


\subsubsection{입력-쿼리 통합 아키텍처}
모델은 입력 센서 토큰과 쿼리 토큰을 단일 시퀀스로 결합하여 처리한다. $N_s$개의 입력 센서와 $N_q$개의 쿼리가 있을 때, 전체 시퀀스 길이는 $N_s+N_q$가 된다. 이러한 통합 처리 방식의 장점은 다음과 같다:
첫째, 쿼리 토큰이 모든 입력 센서의 정보를 직접 참조할 수 있어, 전역적인 환경 패턴을 고려한 예측이 가능하다. 둘째, 각 쿼리가 독립적으로 처리되면서도 동일한 입력 컨텍스트를 공유하여 계산 효율성이 높다. 셋째, 런타임에 쿼리 수와 위치를 자유롭게 변경할 수 있어 유연한 배포가 가능하다.

\subsubsection{쿼리 지정 과정}
실제 예측 시 쿼리 지정은 다음과 같은 단계로 이루어진다:
먼저, 사용자는 예측하고자 하는 위치(3차원 좌표)를 정의한다. 좌표 지정 시에는 실제 물리적 단위를 사용하며, 이는 내부적으로 [0, 1] 범위로 정규화된다. 다음으로, 각 위치에서 예측하고자 하는 센서 타입을 지정한다. 동일한 위치에서 여러 타입을 예측하려면 각 타입별로 별도의 쿼리 토큰을 생성한다. 예를 들어, 한 위치에서 온도와 습도를 동시에 예측하고자 한다면, 두 개의 쿼리 토큰이 생성되며, 각각은 동일한 위치 정보를 공유하지만 다른 타입 임베딩을 갖는다.

쿼리 임베딩 생성 단계에서는 각 쿼리의 위치 정보가 Gaussian Fourier Features를 통해 d차원 벡터로 인코딩되고, 타입 정보는 학습된 임베딩 테이블에서 조회된다. 이 두 임베딩이 합산되어 쿼리 토큰을 형성한다. 입력 토큰과 달리 쿼리 토큰은 측정값에 대한 정보를 포함하지 않으므로, 값 임베딩 없이 위치와 타입 정보만으로 구성된다. 이는 쿼리 토큰이 예측 대상을 지정하는 "질문"의 역할을 수행함을 구조적으로 구현한 것이다.

마지막으로, 여러 시간 윈도우에 대해 연속적인 예측을 수행할 경우, 각 윈도우마다 동일한 쿼리 구성이 반복된다. 24시간 윈도우가 30분씩 이동하면서, 각 윈도우에서 동일한 위치와 타입의 가상 센서값이 예측되어 시간에 따른 연속적인 모니터링이 가능해진다.


\subsubsection{쿼리 토큰의 학습과 추론}
학습 과정에서 쿼리 토큰은 실제 센서 데이터를 타겟으로 사용한다. 즉, 일부 물리 센서를 의도적으로 쿼리로 지정하고, 나머지 센서들로부터 해당 위치의 센서의 값을 예측하도록 학습한다. 이를 통해 모델은 공간적 상관관계와 센서 타입 간 관계를 학습한다.
추론 시에는 물리 센서가 존재하지 않는 임의의 위치를 쿼리로 지정할 수 있다. 모델은 학습된 공간적 패턴을 기반으로, 주변 센서들의 정보를 종합하여 해당 위치의 환경값을 예측한다. 이러한 유연성이 가상 센서 시스템의 핵심 강점이다.

\subsubsection{계층적 어텐션 마스킹 메커니즘}
쿼리 토큰이 올바르게 작동하기 위해서는 정교한 어텐션 마스킹이 필수적이다. 본 연구에서 제안하는 계층적 마스킹 전략은 입력 센서와 쿼리 간의 비대칭적 정보 흐름을 구현한다.
제안된 마스킹 전략은 토큰을 세 가지 타입으로 분류한다: 1) 패딩 토큰, 2) 입력 센서 토큰, 3) 그리고 쿼리 토큰. 각 토큰 타입 간의 어텐션 관계는 다음과 같이 정의된다:
첫째, 입력 센서 토큰들 간에는 완전한 양방향 self-attention이 허용된다. 이를 통해 서로 다른 위치와 타입의 센서들이 상호 정보를 교환하여 전역적인 환경 패턴을 학습할 수 있다. 예를 들어, 축사의 전면부 온도 센서는 후면부 습도 센서의 정보를 참조하여 전체적인 환경 상태를 파악할 수 있다.
둘째, 쿼리 토큰은 모든 입력 센서 토큰을 참조할 수 있지만, 다른 쿼리 토큰들과는 정보를 교환할 수 없다. 이는 각 쿼리가 독립적으로 예측되도록 보장하며, 한 위치의 예측이 다른 위치의 예측에 영향을 주는 것을 방지한다. 이러한 설계는 실제 배포 시나리오를 정확히 반영한다 - 가상 센서는 물리 센서의 정보만을 활용할 수 있으며, 다른 가상 센서의 예측값은 참조할 수 없다.
셋째, 쿼리 토큰은 자기 자신에 대한 self-attention만 허용된다. 이는 어텐션 행렬의 대각선 요소만 활성화되는 것을 의미하며, 각 쿼리가 고유한 표현을 유지하도록 한다.
이러한 마스킹 패턴은 어텐션 점수 행렬에 적절한 마스크 값을 적용하여 구현된다. 허용되지 않는 연결에는 음의 무한대 값을 할당하여, softmax 연산 후 해당 어텐션 가중치가 0이 되도록 한다. 이 마스크는 모든 어텐션 헤드에 동일하게 적용되어 일관된 정보 흐름을 보장한다.
% REVIEW: The hierarchical masking strategy is critical for the "virtual sensing" validity. By preventing query-to-query attention, you ensure parallel queries are independent, which is excellent for scalability. Ensure this independence is emphasized as a feature that allows arbitrary numbers of queries in parallel without interference.


\subsection{데이터 로딩 및 학습 전략}

\subsubsection{동적 센서 분할 메커니즘}
효과적인 모델 학습을 위해 정교한 데이터 로딩 전략을 설계하였다. 학습을 위한 데이터 로더는 각 학습 스텝마다 가용 센서를 동적으로 입력과 타겟으로 분할하여, 모델이 다양한 센서 구성에 대해 강건성을 갖도록 한다.
센서 타입 기반 랜덤화 전략은 다음과 같이 작동한다. 먼저, 네 가지 센서 타입(온도, 습도, CO2, 암모니아) 중 사용할 센서를 1개에서 4개까지 랜덤하게 선택한다. 선택된 각 타입에 대해, 해당 타입의 센서들을 입력과 타겟으로 분할한다. 이때 각 타입별로 최소 1개의 센서는 입력으로 사용되도록 보장하여, 모델이 해당 타입의 기본 패턴을 학습할 수 있도록 한다. 나머지 센서들은 타겟으로 설정되어 예측 대상이 된다.
이러한 접근법의 핵심 이점은 모델이 특정 센서 조합에 과적합되는 것을 방지한다는 점이다. 예를 들어, 어떤 배치에서는 온도 센서 3개와 습도 센서 2개가 입력으로 사용되어 습도 센서를 예측하고, 다른 배치에서는 모든 타입의 센서가 혼재된 입력으로 암모니아 센서를 예측할 수 있다. 이는 실제 배포 환경에서 발생할 수 있는 다양한 센서 고장 시나리오에 대한 모델의 적응력을 향상시킨다.

\subsubsection{결정론적 샘플링과 재현성}
모델 평가의 일관성과 재현성을 보장하기 위해 결정론적 샘플링 옵션을 제공한다. 이 모드에서는 데이터셋 인덱스를 기반으로 고유한 랜덤 시드를 생성하여, 동일한 인덱스는 항상 동일한 센서 분할을 생성한다. 반면, 학습 시에는 비결정론적 모드를 사용하여 매 에폭마다 다른 센서 조합을 경험하도록 함으로써 학습 데이터의 다양성을 극대화한다.

\subsubsection{Self-Reference Training 메커니즘}
본 연구에서 도입한 Self-Reference Training 메커니즘은 모델의 자기 참조 능력을 향상시키기 위한 필수적 구성요소이다. 실험 과정에서 Self-Reference Training 없이 학습된 모델이 특정 상황에서 비최적 동작을 보이는 것이 관찰되었다. 구체적으로, 예측 대상 센서와 동일한 위치 및 타입의 센서가 입력에 포함된 경우에도, 모델이 해당 값을 정확히 재현하지 못하는 현상이 발생하였다.
이론적으로 동일한 위치와 타입의 센서가 입력에 존재한다면, 최적의 예측은 해당 값을 그대로 출력하는 항등 매핑이다. 그러나 Self-Reference Training이 없는 모델은 이러한 단순한 매핑을 학습하지 못하고, 불필요한 보간을 수행하는 경향을 보였다. 이는 전체적인 예측 성능에는 큰 영향을 주지 않지만, 모델의 논리적 일관성 측면에서 문제가 된다.
이 현상의 원인은 다음과 같이 분석된다. 첫째, 쿼리 토큰과 입력 토큰의 구조적 차이가 존재한다. 쿼리 토큰은 값 임베딩이 영벡터로 초기화되는 반면, 입력 토큰은 실제 측정값의 임베딩을 포함한다. 이러한 초기 표현의 차이로 인해 모델이 동일 센서 간 대응 관계를 학습하는 데 어려움을 겪는다.
둘째, 학습 데이터의 분포 특성이 영향을 미친다. 대부분의 학습 샘플에서 쿼리와 정확히 일치하는 입력 센서가 존재하지 않으므로, 모델은 주로 공간 보간 전략을 학습한다. 이로 인해 항등 매핑이 가능한 경우에도 보간을 시도하는 편향이 발생한다.
셋째, 어텐션 메커니즘의 최적화 과정에서 복잡도가 증가한다. 모델이 다양한 센서 조합과 위치 관계를 학습하면서, 단순한 항등 매핑보다 복잡한 패턴을 우선시하는 경향이 나타난다.
Self-Reference Training은 타겟 센서 중 일부를 확률적으로 입력에 포함시킴으로써 이 문제를 해결한다. 실험에서는 0.435의 고정 확률을 사용하였다. 이를 통해 모델은 두 가지 시나리오를 균형있게 학습한다: 동일 센서가 입력에 있을 때의 직접 참조와 없을 때의 공간 보간. 
중요한 설계 결정은 Self-Reference Training된 센서를 타겟에서 제거하지 않는 것이다. 센서가 입력과 타겟 모두에 존재함으로써, 모델은 자기 참조를 명시적으로 학습하게 된다. 이는 일반적인 Self-Reference Training과 구별되는 본 연구의 특징이다.
고정 확률을 사용하는 것은 이 메커니즘이 학습 안정화가 아닌 구조적 문제 해결을 목적으로 하기 때문이다. Curriculum learning과 달리, 전체 학습 과정에서 일관된 비율의 자기 참조 샘플을 제공하는 것이 균형잡힌 학습에 유리하다.
결과적으로 Self-Reference Training을 적용한 모델은 동일 센서가 입력에 있을 경우 높은 정확도의 항등 매핑을 수행하며, 동시에 보간이 필요한 경우에도 우수한 성능을 유지한다. 이는 모델이 상황에 따라 적절한 예측 전략을 선택할 수 있음을 시사한다.
% REVIEW: Terminology Check: "Self-Reference Training" typically refers to using ground-truth *previous* tokens as input in autoregressive sequence generation (RNNs). Here, you are describing a data augmentation or "Identity Mapping Enforcement" strategy where the target is also provided as input. While the mechanism is similar (using ground truth), the term might confuse readers expecting autoregressive behavior. Consider using terms like "Input Masking Augmentation" or "Self-Reference Training" unless "Self-Reference Training" is standard in this specific niche. Also, the necessity of this mechanism suggests the model has a strong bias towards smoothing/interpolation—worth noting why it fails identity mapping (likely due to the strong prior of GFF-based smoothing).


\subsubsection{모델의 해석가능성}
제안된 모델의 중요한 장점 중 하나는 높은 해석가능성이다. 커스텀 Transformer 인코더를 통해 각 레이어의 어텐션 가중치를 추출할 수 있으며, 이를 분석하여 모델의 예측 근거를 이해할 수 있다. 어텐션 가중치는 각 쿼리 위치가 어떤 입력 센서를 주로 참조하는지 보여주며, 이는 센서 중요도 분석과 최적 센서 배치 결정에 활용될 수 있다.
또한, 학습된 타입 임베딩을 분석하여 센서 타입 간의 관계를 이해할 수 있다. 유사한 임베딩 벡터를 가진 센서 타입들은 모델이 유사한 패턴을 가진 것으로 학습했음을 의미하며, 이는 도메인 지식과의 일치성을 검증하는 데 활용될 수 있다.
이러한 해석가능성은 단순한 블랙박스 예측을 넘어, 모델의 의사결정 과정을 이해하고 신뢰할 수 있는 가상 센서 시스템 구축을 가능하게 한다. 이는 특히 안전이 중요한 축사 환경 모니터링 응용에서 필수적인 특성이다.

\subsubsection{모델 구조 및 하이퍼파라미터}

제안된 모델의 하이퍼파라미터는 광범위한 실험을 통해 최적화되었다. 최종 선택된 하이퍼파라미터는 다음과 같다:
모델 아키텍처 파라미터로는 히든 차원($d_{\text{model}}$)이 768로 설정되어 고차원 표현 학습이 가능하도록 하였다. Transformer 인코더는 3개 레이어($N_{\text{layer}}=3$)로 구성되며, 각 레이어는 8개의 어텐션 헤드($N_{\text{head}}=8$)를 사용한다. 이는 $d_{\text{model}}$을 헤드 수로 나눈 값이 96차원이 되어, 각 헤드가 충분한 표현력을 갖도록 한다. Feed-forward network의 히든 차원은 3072($4 \times d_{\text{model}}$)으로 설정되었다. 드롭아웃 확률($p_{\text{dropout}}$)은 0.27로 설정되어 적절한 regularization을 제공한다.


MLP 구조는 입력과 출력 단계에서 다르게 구성된다. Value MLP는 3개 레이어($N_{VP}=3$)로 구성되어 48개 타임스텝의 시계열 데이터를 $d_{\text{model}}$ 차원으로 효과적으로 변환한다. 반면, Output MLP는 단일 레이어($N_{OL}=2$)로 구성하여 과적합을 방지하면서도 필요한 역변환을 수행한다.
학습 관련 파라미터로는 배치 크기($B_{\text{train}}$)를 128로 설정하여 메모리 효율성과 학습 안정성의 균형을 맞추었다. Self-Reference Training 확률($p_{\text{self}}$)은 0.435로 고정되어, 전체 학습 과정에서 일관되게 자기 참조 학습을 유도한다. 시계열 윈도우 크기는 48 타임스텝(24시간)으로 설정되어 일중 변화 패턴을 포착한다.

Gaussian Fourier Features의 표준편차도 하이퍼 파라미터로 최적화 되었으나 우리의 실험 결과에서는 너무 큰 $\sigma$ (예, $\sigma=10$)값을 사용하면 공간 interpolation이 작동하지 않는 것을 확인했다.  선택된 하이퍼파라미터는 Table. \ref{tab:hyperparams}에 정리하였다.


\begin{table}[t]
    \centering
    % 행 간격을 살짝 넓혀서 시원하게 만듭니다 (기본값 1.0)
    \renewcommand{\arraystretch}{1.0} 
    
    \resizebox{\columnwidth}{!}{%
    \begin{tabular}{l l c} %% 정렬 변경: 좌측(l), 좌측(l), 중앙(c)
        \toprule
        \textbf{Symbol} & \textbf{Description} & \textbf{Value} \\
        \midrule
        $B_{\text{train}}$ & Training batch size & 128 \\
        $d_{\text{model}}$ & Hidden size of model & 768 \\
        $N_{\text{layer}}$ & Number of layers in encoder& 3 \\
        $N_{\text{head}}$ & Number of MHA heads & 8 \\
        $p_{\text{dropout}}$ & Dropout rate & 0.27 \\
        $N_{VP}$ & Number of layers in value projector MLP& 3 \\
        $N_{OL}$ & Number of layers in output head MLP& 2 \\
        $p_{\text{self}}$ & Self-Reference Training prob. & 0.435 \\
        lr & Adam learning rate & 0.00013 \\
        $\sigma$ & Sigma & 1.37 \\
        \bottomrule
    \end{tabular}%
    }
    \caption{Selected Hyperparameters}
    \label{tab:hyperparams}
\end{table}

\subsubsection{학습 최적화 및 손실 함수}
모델 학습에는 Mean Squared Error (MSE)를 손실 함수로 사용한다. MSE는 예측값과 실제값 간 차이의 제곱을 평균하여 계산되며, 큰 오차에 대해 더 큰 페널티를 부과한다. 이는 가상 센서의 정확도가 중요한 본 응용에서 큰 예측 오차를 효과적으로 줄이는 데 적합하다.
최적화 알고리즘으로는 Adam optimizer를 사용하며, 학습률은 0.00013으로 설정되었다. 이는 하이퍼파라미터 최적화를 통해 선택된 값으로, 안정적인 수렴과 빠른 학습 속도를 보장한다.
학습은 스텝 기반으로 진행되며, 매 100 스텝마다 검증을 수행한다. 이는 에폭 기반 학습보다 세밀한 모니터링과 early stopping을 가능하게 한다. Early stopping은 검증 손실이 10회 연속 개선되지 않을 경우 발동되며, 이는 patience=10 스텝으로 설정되었다.

\subsection{배치 최적화 과정}
이 연구에서는 센서의 최적 배치를 각 타입별로 분류해서 수행
테스트 구간에서의 센서의 수는 온습도 센서 각각 8개, 이산화탄소 센서 9개, 암모니아 센서 6개
모든 경우의 수를 계산할 수 있기 때문에 특별한 최적화 알고리즘을 사용하지 않고 brute-force 방법을 사용함. 
% REVIEW: The term "brute-force" suggests an exhaustive search. You must justify why the search space is small enough for this to be feasible. With 8+8+9+6 sensors, the combination space could be explosive ($C(n, k)$) unless the candidate locations are very limited. Provide the size of the search space (e.g., "Total combinations evaluated: X"). If it's truly exhaustive, calculate the complexity. If it's greedy or constrained, describe it accurately.




%% Use \section commands to start a section
\section{Experimental design}

%% Use \subsection commands to start a subsection.
\subsection{aaaa}

%% Use \section commands to start a section
\section{Results and discussion}

%% Use \subsection commands to start a subsection.
\subsection{학습 결과 및 어텐션 분석}
본 섹션에서는 학습된 모델의 시계열 예측 결과와 어텐션 맵 분석을 통합하여 모델의 동작 원리를 다각도로 분석한다. 어텐션 가중치는 모든 인코더 레이어, multi-head attention의 모든 헤드, 그리고 모든 예측 배치에 대하여 평균내어 계산되었다.

첫째, 타겟 센서를 제외한 모든 물리 센서를 입력으로 사용하여 온도 센서 예측을 수행한 결과, 모델은 매우 높은 예측 정확도를 보였다. 그러나 어텐션 맵 분석 결과, 같은 타입(온도) 센서에 대한 어텐션 가중치는 높은 반면, 다른 타입 센서(습도, CO2, 암모니아)에 대한 어텐션은 매우 낮게 나타났다. 이는 모델이 동일 타입 센서의 존재 유무에만 주목하고, 이종 센서 간의 복잡한 물리적 상관관계보다는 동종 센서 간의 직접적인 패턴 유사성을 학습하는 데 집중했기 때문으로 해석된다. 본 논문에서는 온도 센서를 대표 사례로 제시하지만, 이는 모든 센서 타입에 대해 전반적으로 관찰되는 현상임을 명시한다.

\begin{figure*}[htbp]
  \centering
  \includegraphics[width=\textwidth]{temp_all_sensors.pdf}
  \captionsetup{width=\textwidth}
  \caption{온도 센서 예측 결과 및 어텐션 분석. 타겟 센서를 제외한 모든 물리 센서를 사용한 경우. (a) 시계열 예측 결과 (검은 선: ground-truth, 빨간 선: 예측 결과). (b) 어텐션 맵 (원의 크기가 중요도를 나타냄, 온도 센서에 대한 어텐션이 높고 다른 타입 센서에 대한 어텐션은 매우 낮음).}
  \label{fig:temp_all_sensors}
\end{figure*}

둘째, 습도 예측에 대하여 같은 타입(습도) 센서만을 사용하여 예측을 수행한 결과, 성능이 좋게 예측됨을 확인하였다. 어텐션 맵 분석 결과, 예측 위치에서 가장 가까운 습도 센서의 중요도가 가장 높게 나타났으며, 이는 공간적으로 인접한 센서 간의 높은 상관관계를 반영하는 것으로 물리적으로 타당한 결과이다. 이는 모델이 공간적 근접성을 효과적으로 학습하고 있음을 시사한다.

\begin{figure*}[htbp]
  \centering
  \includegraphics[width=\textwidth]{humid_same_sensors.pdf}
  \captionsetup{width=\textwidth}
  \caption{습도 센서 예측 결과 및 어텐션 분석. 같은 타입(습도) 센서만 사용한 경우. (a) 시계열 예측 결과 (검은 선: ground-truth, 파란 선: 예측 결과). (b) 어텐션 맵 (원의 크기가 중요도를 나타냄, 근처 센서의 중요도가 가장 높음).}
  \label{fig:humidity_same_type}
\end{figure*}

셋째, CO2 센서에 대하여 센서 중요도 기반 제거 실험을 수행하였다. 어텐션 가중치를 기반으로 가장 중요한 3개 센서와 가장 안 중요한 3개 센서를 각각 식별한 후, 각각의 경우에 대해 센서를 제거하고 예측 성능을 비교하였다. 실험 결과, 가장 중요한 3개 센서를 제거한 경우 예측 성능이 크게 저하된 반면, 가장 안 중요한 3개 센서를 제거한 경우에는 성능 저하가 상대적으로 적게 나타났다. 이는 어텐션 가중치가 센서의 실제 중요도를 잘 반영하고 있음을 보여주며, 모델의 해석가능성과 신뢰성을 뒷받침한다.

\begin{figure*}[htbp]
  \centering
  \begin{tabular}{cc}
    \fbox{\parbox{0.45\textwidth}{\centering \vspace{2cm} PLACEHOLDER (a) \vspace{2cm}}} &
    \fbox{\parbox{0.45\textwidth}{\centering \vspace{2cm} PLACEHOLDER (b) \vspace{2cm}}} \\
    (a) 가장 중요한 3개 센서 제거 & (b) 가장 안 중요한 3개 센서 제거 \\
  \end{tabular}
  \captionsetup{width=\textwidth}
  \caption{CO2 센서 중요도 기반 제거 실험 결과. 각 subfigure는 시계열 예측 결과(상단)와 어텐션 맵(하단)을 포함. (a) 가장 중요한 3개 센서를 제거한 경우 (성능 저하 큼). (b) 가장 안 중요한 3개 센서를 제거한 경우 (성능 저하 적음).}
  \label{fig:co2_removal}
\end{figure*}

% REVIEW: The integrated analysis of prediction results and attention maps provides a comprehensive view of model behavior. The finding that same-type sensors dominate attention weights, while physically reasonable, suggests limited multi-modal sensor fusion. The CO2 removal experiment effectively validates the interpretability of attention weights, demonstrating that they reflect actual sensor importance.

\subsection{공간 Interpolation 분석}
모델은 물리 센서가 설치된 격자 사이의 임의 위치에 대해서도 예측을 수행하였으며, 이를 통해 연속적인 환경 변수 값의 변화를 관찰할 수 있었다. 비록 해당 위치의 실제 측정값(Ground Truth)이 부재하여 정량적 검증에는 한계가 있으나, 시각적으로 타당한 연속성을 보였다.

한 가지 주목할 점은 위치 임베딩 생성 시 사용되는 $\sigma$ 값의 영향이다. $\sigma$ 값이 너무 클 경우, 미세한 위치 변화에도 임베딩 벡터가 급격하게 변하여 공간적 연속성이 깨지고, 사잇값이 랜덤하게 예측되는 현상이 발견되었다. 이는 고주파 성분이 과도해져 위치 임베딩이 사실상 랜덤 벡터처럼 동작하기 때문으로 분석된다. 또한, 공간 내 센서 수가 제한적인 점도 이러한 현상에 영향을 미친 것으로 보인다.

\begin{figure}[htbp]
  \centering
  \fbox{\parbox{0.45\columnwidth}{\centering \vspace{1.5cm} PLACEHOLDER (a) \vspace{1.5cm}}}
  \fbox{\parbox{0.45\columnwidth}{\centering \vspace{1.5cm} PLACEHOLDER (b) \vspace{1.5cm}}}
  \caption{공간 인터폴레이션 결과 (특정 타임스텝). (a) 온도, (b) 이산화탄소.}
  \label{fig:interp_results}
\end{figure}

% REVIEW: The discussion on sigma sensitivity is insightful. Providing a specific range of stable sigma values based on your experiments would make this finding more actionable for future researchers.

\subsection{최적 센서 배치 (동일 타입)}
각 센서 타입별로 물리 센서의 개수에 따른 최적의 배치 전략을 분석하였다. 본 실험에서는 다른 타입의 센서 영향을 배제하고 동일 타입 센서만을 고려하였으며, 최적 배치의 기준으로는 가상 센서들의 RMSE 평균값이 최소가 되는 조합을 선정하였다. 이때 MAE 대신 RMSE를 사용하여 큰 예측 오차에 더 큰 페널티를 부여하였다.

\subsubsection{최적 물리 센서 배치 결과}
분석 결과, 물리 센서를 1개만 사용할 경우 모든 타입에서 직관적으로 공간의 중앙에 위치한 센서가 최적 위치로 선정되었다. 단, 암모니아 센서의 경우 가장 아래쪽 센서가 선택되었는데, 이는 해당 환경의 특수성(예: 가스 비중 등)에 기인한 것으로 추정되며 추가 분석이 필요하다.

2개 이상의 물리 센서를 사용하는 경우에는 중앙 센서가 포함되는 빈도가 낮았다. 이는 주변 센서들의 값을 통해 중앙부의 예측이 비교적 용이하기 때문으로 해석된다. 이산화탄소를 제외하면, 일반적으로 $n$개일 때의 최적 배치는 $n-1$개의 최적 배치에 새로운 센서를 하나 추가하는 형태를 띠는 경향을 보였다.

\begin{figure*}[htbp]
  \centering
  \begin{tabular}{cc}
    \fbox{\parbox{0.9\columnwidth}{\centering \vspace{2cm} PLACEHOLDER (a) \vspace{2cm}}} &
    \fbox{\parbox{0.9\columnwidth}{\centering \vspace{2cm} PLACEHOLDER (b) \vspace{2cm}}} \\
    (a) Temperature & (b) Humidity \\
    \fbox{\parbox{0.9\columnwidth}{\centering \vspace{2cm} PLACEHOLDER (c) \vspace{2cm}}} &
    \fbox{\parbox{0.9\columnwidth}{\centering \vspace{2cm} PLACEHOLDER (d) \vspace{2cm}}} \\
    (c) CO2 & (d) NH3 \\
  \end{tabular}
  \caption{물리 센서 개수별 최적 배치 시각화. (a) 온도, (b) 습도, (c) 이산화탄소, (d) 암모니아.}
  \label{fig:opt_placement}
\end{figure*}

\subsubsection{물리 센서 개수별 성능 트렌드}
물리 센서의 개수와 배치에 따른 성능 변화 추이를 분석하였다. 이때 성능 평가는 예측 대상이 되는 모든 가상 센서 위치에서의 RMSE 값의 산술 평균을 기준으로 하였다.

물리 센서가 특정 위치에 배치되어야 하는 이유는 크게 두 가지 관점에서 고려해볼 수 있다. 첫째는 해당 센서의 값이 다른 위치의 센서 값을 예측하는 데 중요한 정보원(Information Source)으로 작용하여 전체적인 예측 정확도를 높이는 경우이다. 둘째는 해당 위치의 환경 변수 자체가 주변 센서들로부터 예측하기 어려워, 직접적인 측정이 불가피한 경우이다. 이론적으로는 이 두 가지 목적이 구분될 수 있으나, 가상 센싱 시스템의 실질적인 운영 목표는 개별 센서의 역할 구분보다는 공간 전체의 평균적인 예측 오차(RMSE)를 최소화하는 것이다. 따라서 본 분석에서는 구분 없이 전체 RMSE 감소를 최우선 지표로 삼았다.

% REVIEW: The distinction between "informational" sensors and "hard-to-predict" targets is a sophisticated insight. It adds depth to the placement strategy discussion. Prioritizing global RMSE aligns well with the practical goal of system optimization.

분석 결과, 동일한 모델과 동일한 개수의 센서를 사용하더라도, 센서의 공간적 배치에 따라 예측 성능이 극적으로 달라짐을 확인하였다. 특히 주목할 점은, 배치가 최적화되지 않은 경우 물리 센서의 개수를 늘리더라도 평균적인 성능이 오히려 저하되는 현상이 빈번하게 발생했다는 것이다. 이는 부적절한 위치의 센서 데이터가 모델에게 유용한 정보를 제공하기보다는 오히려 노이즈로 작용하거나, 중복된 정보를 제공하여 자원의 비효율을 초래할 수 있음을 시사한다.

반면, 제안된 최적 배치 전략을 따를 경우, 물리 센서의 수가 증가함에 따라 예측 성능(RMSE)이 일관되게 향상되는 경향을 보였다. 이는 "데이터는 많을수록 좋다"는 일반적인 통념이 항상 성립하는 것은 아니며, "어떤 데이터를 수집하느냐"가 훨씬 더 중요함을 실증적으로 보여준다. 결론적으로, 사용자는 가용 예산과 요구되는 정확도 수준을 고려하여 물리 센서의 개수를 유연하게 선택하되, 반드시 최적의 위치 선정이 선행되어야 함을 알 수 있다.

% REVIEW: The observation that Adding more sensors can *degrade* performance if poorly placed is a counter-intuitive and strong finding. It strongly validates the necessity of the proposed placement algorithm. The conclusion calling for "quality over quantity" in data collection is compelling.

\begin{figure*}[htbp]
  \centering
  \begin{tabular}{cc}
    \fbox{\parbox{0.9\columnwidth}{\centering \vspace{2cm} PLACEHOLDER (a) \vspace{2cm}}} &
    \fbox{\parbox{0.9\columnwidth}{\centering \vspace{2cm} PLACEHOLDER (b) \vspace{2cm}}} \\
    (a) Temperature & (b) Humidity \\
    \fbox{\parbox{0.9\columnwidth}{\centering \vspace{2cm} PLACEHOLDER (c) \vspace{2cm}}} &
    \fbox{\parbox{0.9\columnwidth}{\centering \vspace{2cm} PLACEHOLDER (d) \vspace{2cm}}} \\
    (c) CO2 & (d) NH3 \\
  \end{tabular}
  \caption{물리 센서 개수 및 배치에 따른 가상 센서 예측 RMSE 평균 (Red: 최적 배치, Gray: 그 외). (a) 온도, (b) 습도, (c) 이산화탄소, (d) 암모니아.}
  \label{fig:perf_trends}
\end{figure*}

% REVIEW: The observation that non-optimal placement of additional sensors can degrade average performance is a key contribution. It contradicts the common assumption that "more data is better" and highlights the importance of the proposed optimization framework. Ensure this point is highlighted in the conclusion as well.



%% Use \section commands to start a section
\section{Conclusions}


%% Use \subsection commands to start a subsection.
\subsection{aaaa}



% \end{linenumbers}

\section*{Acknowledgments}
This work was supported by the National Research Foundation of Korea (NRF) grant (No. RS-2023-00277318 and RS-2025-00512551) funded by the Korean government (Ministry of Science and ICT).



상호 참조하기 -> Fig. \ref{fig:fig1}
%% The Appendices part is started with the command \appendix;
%% appendix sections are then done as normal sections
\appendix
\renewcommand{\figurename}{Fig} % "Figure"를 "Fig"로 변경
\renewcommand{\thefigure}{\Alph{section}\arabic{figure}} % "Appendix B" 대신 "B"만 가져옴
\renewcommand{\thetable}{\Alph{section}\arabic{table}} % 표 번호도 동일하게 적용

%% Appnedix A
\section{Appendix A. aaaa}
\label{appendA}
\setcounter{figure}{0}
\setcounter{table}{0}


Appendix text.


%% Appendix B 
\section{Appendix B. bbbbb}
\label{appendB}
\setcounter{figure}{0}
\setcounter{table}{0}

Appendix B text 

\begin{figure}[htbp]%% placement specifier
%% Use \includegraphics command to insert graphic files. Place graphics files in 
%% working directory.
\centering%% For centre alignment of image.
\includegraphics[width=\columnwidth]{pptl-diagram.png}
%% Use \caption command for figure caption and label.
\caption{{컬럼 너비 피겨 (옵션 \texttt{width=\textbackslash columnwidth} 사용)}}
\label{fig:appendix-fig1}
%% https://en.wikibooks.org/wiki/LaTeX/Importing_Graphics#Importing_external_graphics
\end{figure}





%% Supplementary material 섹션
\section*{Supplementary material}
Interactive plot for time-series pattern exploration: An interactive HTML extension of Fig. \ref{fig:humidity_same_type}(a) is provided with two key features

%% Code and data availability 섹션
\section*{Code and data availability}

Source code and data used in this study are available at \url{https://github.com/bet-lab/CEiT}.


%% If you have bib database file and want bibtex to generate the
%% bibitems, please use
%%

% \bibliographystyle{elsarticle-num} 
\bibliographystyle{elsarticle-harv} 
\bibliography{references}

%%  \bibliography{<your bibdatabase>}

%% else use the following coding to input the bibitems directly in the
%% TeX file.

%% Refer following link for more details about bibliography and citations.
%% https://en.wikibooks.org/wiki/LaTeX/Bibliography_Management



%% Examples 섹션은 필요시 주석 해제하여 사용하세요
% \section{Examples}
% \subsection{Figures}
% \begin{figure}[htbp]
%   \centering
%   \includegraphics[width=\linewidth]{figure_name} % 확장자(.eps, .pdf) 생략 가능
%   \caption{Description of the figure.}
%   \label{fig:my_figure}
% \end{figure}



\end{document}
